\section{Uniform power decay and the gradient estimate}\label{sec:reduction}

We now prove the estimate in Theorem~\ref{thm:main}, conditional on a Liouville Theorem whose proof occupies the remaining sections. The argument works with all normalized solutions at once. We first obtain a common modulus of affine approximation, then use anisotropic compactness to force a power rate. Fits at this rate give the H\"older estimate for the gradient. The scaling and limiting equation also occur in Evans and Savin \cite[Sections~2.3--2.4]{EvansSavin2008}.

\begin{theorem}[Liouville Theorem]\label{thm:liouville}
Let $m\ge1$ and $0\le C_0<\infty$. Suppose that $v\in C(\R^m\times\R)$ is a viscosity solution of
\begin{equation}\label{eq:limit-pde}
  F[v]:=D^2v[(\nabla_xv,1)]=0.
\end{equation}
Suppose, in addition, that for every $z=(x,t)\in\R^m\times\R$ and every $R>0$ there are $p\in\R^m$ and $a\in\R$ such that
\begin{equation}\label{eq:limit-fit}
  \left|v(x+h,t+\ell)-v(x,t)-p\cdot h
       -\left(a-\frac{|p|^2}{2}\right)\ell\right|\le C_0R
  \quad\text{if }\quad |\ell|\le R,\quad |h-p\ell|\le R.
\end{equation}
Then $v$ is affine.
\end{theorem}

We prove Theorem~\ref{thm:liouville} in Sections~\ref{sec:limiting-tools}--\ref{sec:rigidity}. For clarity, a pair $(p,a)$ satisfying \eqref{eq:limit-fit} will be called a \emph{fit at $z$ of radius $R$}. Its domain is the sheared box
\[
  Q_R(z;p)=\{(x+h,t+\ell): |\ell|\le R,\ |h-p\ell|\le R\}.
\]
The ordinary time coefficient of the fitted affine function is $q=a-|p|^2/2$. We use $(p,a)$ because a sheared change of frame acts by subtracting the frame parameters, as shown in Section~\ref{sec:limiting-tools}.

\subsection{A flatness modulus for the normalized family}

Fix $d\ge3$, put $n=d+1$, and write $B_r^k$ for the ball of radius $r$ centered at zero in $\R^k$. Let $\mathcal U$ be the family of all functions
\[
  U(y,b)=u(y)+b\qquad ((y,b)\in B_1^d\times\R),
\]
where $u$ is a continuous infinity-harmonic function with
$u(0)=0$ and $\osc_{B_1^d}u\le1$. By Lemma~\ref{lem:linear-lift}, each $U$ is infinity-harmonic. The family is uniformly bounded and has a common Lipschitz constant $L$ on $\overline B_{7/8}^n$. Indeed, $|u|\le1$, and the oscillations of the lifts are uniformly bounded on $\overline B_{15/16}^n$. At every center in $\overline B_{7/8}^n$, Lemma~\ref{lem:cone-slopes} bounds both slopes of radius $1/16$ by that common oscillation divided by $1/16$. This gives a common Lipschitz bound for nearby pairs; subdividing a segment gives the bound for all pairs in the ball. All constants in this section, unless indexed by additional parameters, depend only on $d$.

Set
\[
  G=B_{3/4}^n,\qquad D(z)=\operatorname{dist}(z,\partial G),
\]
and denote the added unit direction by $e_b$. For $P\ne0$, $0<\lambda<1$, and $r>0$, put $N_P=P/|P|$ and define the closed cylinder
\[
  Z_\lambda(z,r;P)
   =\{z+h: |h\cdot N_P|\le r,\quad
               |h-(h\cdot N_P)N_P|\le\sqrt\lambda r\}.
\]
For $U\in\mathcal U$, $z\in G$, and $0<r\le D(z)/8$, call $\lambda$ \emph{good for $U$ at $(z,r)$} if some $P\ne0$ satisfies
\begin{equation}\label{eq:physical-fit}
  |U(Z)-U(z)-P\cdot(Z-z)|\le\lambda r
       \qquad (Z\in Z_\lambda(z,r;P)).
\end{equation}
We call this a cylinder fit. Its cylinder is contained in $B_{2r}(z)\subset G$. Let $\beta_U(z,r)$ be the supremum of the parameters in $(0,1)$ that are not good, with value zero if there are none, and set
\begin{equation}\label{eq:flatness-function}
  b(s)=\sup\{\beta_U(z,r):U\in\mathcal U,\ z\in G,
                        \ 0<r\le sD(z)\},\qquad 0<s\le1/8.
\end{equation}
Thus $0\le b\le1$ and $b$ is nondecreasing. Goodness itself need not be monotone in $\lambda$. The supremum definition ensures that every parameter in $(\beta_U(z,r),1)$ is good.

\begin{lemma}[Uniform qualitative flatness]\label{lem:family-flatness}
The family modulus in \eqref{eq:flatness-function} satisfies
\begin{equation}\label{eq:family-flatness}
  b(s)\longrightarrow0\qquad(s\downarrow0).
\end{equation}
Every cylinder fit satisfies
\begin{equation}\label{eq:physical-slope-bounds}
  |P|\le L+\sqrt\lambda,\qquad
  |P\cdot e_b-1|\le\sqrt\lambda.
\end{equation}
\end{lemma}

\begin{proof}
The cylinder contains $\overline B_{\sqrt\lambda r}(z)$. Evaluating the fit on this ball and using the Lipschitz bound gives the first inequality in \eqref{eq:physical-slope-bounds}. The exact increment of $U$ in direction $e_b$ equals the increment length, which gives the second.

Suppose that \eqref{eq:family-flatness} fails. There are $s_i\to0$, functions $U_i\in\mathcal U$, centers $z_i\in G$, radii $0<r_i\le s_iD(z_i)$, and parameters $\lambda_i$ that are not good at these triples, with $\lambda_i\ge\delta>0$. Passing to a subsequence, the common bounds give
\[
  z_i\longrightarrow z_*\in\overline G,\qquad
  U_i\longrightarrow U_\circ
       \quad\hbox{uniformly on }\overline B_{7/8}^n.
\]
The limit $U_\circ$ is infinity-harmonic on $B_{7/8}^n$ by the usual viscosity stability argument \cite[Section~6]{CIL1992}: add a positive fourth-order term to an upper test to make its contact strict, take nearby interior maxima for the approximating functions minus that test, and pass their inequalities to the limit. The added term leaves the first two derivatives at contact unchanged. A negative fourth-order term and minima treat lower tests. This argument also applies to the varying continuous jet operators below, since they converge uniformly on compact jet sets. Choose scales $R_i\downarrow0$ so slowly that
\begin{equation}\label{eq:slow-scales}
  r_i+|z_i-z_*|+
    \|U_i-U_\circ\|_{L^\infty(\overline B_{7/8}^n)}=o(R_i).
\end{equation}
Such a decreasing sequence can be obtained after passage to a subsequence, since all three terms tend to zero. The point $z_*$ is interior to $B_{7/8}^n$. Proposition~\ref{prop:affine-blowups}, applied to the one function $U_\circ$ at $z_*$, gives a further subsequence on which its $R_i$ blow-ups converge locally uniformly to a linear function. By \eqref{eq:slow-scales}, the functions
\[
  \widehat U_i(X)=
    \frac{U_i(z_*+R_iX)-U_i(z_*)}{R_i}
\]
converge to that same linear function uniformly on $B_1^n$. They are infinity-harmonic and uniformly Lipschitz there. Proposition~\ref{prop:three-point} therefore gives numbers $\varepsilon_i\to0$ such that
\begin{equation}\label{eq:slow-scale-midpoint}
  |\widehat U_i(X)+\widehat U_i(Y)
       -2\widehat U_i((X+Y)/2)|\le\varepsilon_i|X-Y|
       \qquad (X,Y\in B_{1/4}^n).
\end{equation}

Return now to the smaller radii $r_i$. The functions
\[
  W_i(H)=\frac{U_i(z_i+r_iH)-U_i(z_i)}{r_i}
\]
are uniformly bounded and Lipschitz on $\overline B_2^n$, so a further subsequence converges uniformly there to $W$. The slow-scale choice ensures that $z_i+r_i\overline B_2^n\subset B_{R_i/4}(z_*)$ for all large $i$. Substituting
\[
  X=\frac{z_i-z_*+r_iH}{R_i},\qquad
  Y=\frac{z_i-z_*+r_iK}{R_i}
\]
into \eqref{eq:slow-scale-midpoint} shows that the midpoint defect of $W_i$ is at most $\varepsilon_i|H-K|$. Thus $W$ satisfies the midpoint identity on $\overline B_2^n$. Continuity and dyadic subdivision on segments make $W$ affine on this ball. Since $W(0)=0$ and $W(e_b)=1$, it has the form $W(H)=P\cdot H$ with $P\cdot e_b=1$.

Uniform convergence now gives, for all large $i$,
\[
  \sup_{|H|\le2}|U_i(z_i+r_iH)-U_i(z_i)-r_iP\cdot H|
       <\delta r_i.
\]
The nonzero slope $P$ therefore gives a fit with error at most $\lambda_i r_i$ on $Z_{\lambda_i}(z_i,r_i;P)\subset B_{2r_i}(z_i)$. This contradicts the choice of $\lambda_i$ and proves \eqref{eq:family-flatness}.
\end{proof}

The larger scales $R_i$ allow a varying sequence of solutions to be compared with a blow-up of one limit function. The three-point estimate then transfers flatness back to the possibly much smaller scales $r_i$. No uniform differentiability theorem for the original family is used.

\subsection{Selection of scales}

Our next objective is a power bound for the single modulus $b$, which already takes the supremum over all normalized solutions.

\begin{proposition}[Uniform power decay]\label{prop:power-decay}
Assume Theorem~\ref{thm:liouville}. There are $B\in(0,\infty)$, $\gamma\in(0,1]$, and $s_0\in(0,1/8]$, depending only on $d$, such that
\begin{equation}\label{eq:power-decay}
  b(s)\le Bs^\gamma\qquad(0<s\le s_0).
\end{equation}
\end{proposition}

We prove this proposition by contradiction. Its failure means that $b(s)/s^\gamma$ is unbounded near zero for every $\gamma>0$. The next lemma selects scales where this failure occurs, while obtaining a slightly enlarged good parameter uniformly at nearby centers and radii.

\begin{lemma}\label{lem:bad-scales}
Suppose that \eqref{eq:power-decay} fails for every positive exponent. There are $\gamma_i\downarrow0$, $s_i\downarrow0$, functions $U_i\in\mathcal U$, centers $z_i\in G$, radii $0<r_i\le s_iD(z_i)$, and parameters $\lambda_i\downarrow0$ such that $\lambda_i$ is not good for $U_i$ at $(z_i,r_i)$ and the following holds. Set $\Lambda=32$. For every fixed $M\ge2$, for all sufficiently large $i$, every center $z$ with $|z-z_i|\le Mr_i$ and radius $0<r\le Mr_i$ are admissible, and $\Lambda\lambda_i$ is good for $U_i$ at $(z,r)$.
\end{lemma}

\begin{proof}
Choose any sequence $\gamma_i\downarrow0$. The function $b(s)s^{-\gamma_i}$ is bounded on intervals separated from zero and is unbounded near zero. Taking a factor-two approximate maximum on a sufficiently long truncated interval $[h_i,1/8]$ gives a point $s_i$ such that $s_i\to0$ and
\begin{equation}\label{eq:running-maximum}
  \frac{b(s)}{s^{\gamma_i}}
     \le2\frac{b(s_i)}{s_i^{\gamma_i}}\qquad(s_i\le s\le1/8).
\end{equation}
For example, make the approximate maximum larger than the supremum on $[1/i,1/8]$ to ensure $s_i<1/i$. Pass to a subsequence to make $s_i$ decreasing. We have $b(s_i)>0$. Choose $U_i\in\mathcal U$, $z_i\in G$, and $r_i\le s_iD(z_i)$ with $\beta_{U_i}(z_i,r_i)>b(s_i)/2$, and then choose a parameter $\lambda_i>\beta_{U_i}(z_i,r_i)/2$ that is not good for $U_i$ at this center and radius. This uses only approximate suprema and gives
\begin{equation}\label{eq:bad-parameter}
  b(s_i)/4<\lambda_i\le b(s_i).
\end{equation}
In particular $\lambda_i\to0$; pass to a further subsequence if needed to make it decreasing.

Fix $M\ge2$. For large $i$, we have $Mr_i\le Ms_iD(z_i)<D(z_i)$, so every $z$ with $|z-z_i|\le Mr_i$ belongs to $G$. The distance to the boundary then satisfies
\[
  D(z)\ge D(z_i)-Mr_i\ge\tfrac12D(z_i)>0.
\]
For $r\le Mr_i$ we have $r/D(z)\le2Ms_i\le1/8$, so the center and radius are admissible. By \eqref{eq:running-maximum} and \eqref{eq:bad-parameter}, uniformly over all these choices,
\[
  \beta_{U_i}(z,r)\le b(2Ms_i)
     \le2b(s_i)(2M)^{\gamma_i}
     \le8\lambda_i(2M)^{\gamma_i}
     <32\lambda_i<1
\]
for all sufficiently large $i$. Here $(2M)^{\gamma_i}\to1$ for each fixed $M$, so the same $\Lambda=32$ works for every fixed $M$; only the threshold in $i$ depends on $M$. The parameter $\Lambda\lambda_i$ is strictly greater than every bad parameter and is therefore good.
\end{proof}

We isolate the geometry used to compare fitted slopes. It is a statement about affine fits and does not use the equation. When $\theta^2=\Lambda\lambda_i$, the rescaling below divides transverse slope errors by $\sqrt{\lambda_i}$ and longitudinal errors by $\lambda_i$. The $O(\theta)$ direction estimate controls the normalized transverse coefficient, while the $O(\theta^2)$ norm estimate, together with the quadratic angular error, controls the normalized longitudinal coefficient.

\begin{lemma}[Comparison of cylinder fits]\label{lem:cylinder-comparison}
Fix $0<c_*\le C_*<\infty$. Suppose two fits for the same function, of the form \eqref{eq:physical-fit}, with parameter $\theta^2\in(0,1)$, have slopes $P,P'$ satisfying $c_*\le |P|,|P'|\le C_*$. Write $n=P/|P|$ and $n'=P'/|P'|$. If their centers agree and their radii are $H$ and $H'$ with $H\le H'\le2H$, then
\begin{equation}\label{eq:cylinder-comparison}
  |P-P'|+|n-n'|\le C\theta,\qquad
  \bigl||P|-|P'|\bigr|\le C\theta^2,
\end{equation}
where $C$ depends only on $c_*,C_*$. The same conclusion holds for equal radii $H$ at centers $z,z'$ if $z'$ lies in the first cylinder with axial and transverse offsets at most $H/4$ and $\theta H/4$, respectively.
\end{lemma}

\begin{proof}
In the common-center case both cylinders contain $B_{\theta H}(z)$. Evaluation of the difference of the fits on this ball gives $|P-P'|\le3\theta$. Normalization of the vectors, using the lower norm bound, gives $|n-n'|\le C\theta$. Choose a fixed small $c>0$. The point $z+cHn$ belongs to both cylinders: relative to $n'$ its transverse displacement is at most $cH|n-n'|\le Cc\theta H$, and its axial displacement is at most $cH$. Taking $c$ sufficiently small, the fit errors at that point give
\[
  |(P-P')\cdot n|\le C\theta^2.
\]
Since $1-n'\cdot n=|n-n'|^2/2$,
\[
  \bigl||P|-|P'|\bigr|
   \le |(P-P')\cdot n|+|P'|(1-n'\cdot n)\le C\theta^2.
\]

For displaced centers, a ball $B_{c\theta H}(z')$ with fixed small $c$ belongs to both cylinders by the quarter-radius margins. Subtracting the first fit's value at $z'$ makes it an estimate for increments from $z'$ with error $2\theta^2H$; the second fit has error $\theta^2H$. Evaluation on the common ball again gives the vector and direction bounds. The segment $z'+s n'$ for $|s|\le c'H$ lies in both cylinders when $c'>0$ is sufficiently small: the first cylinder's initial offsets use a quarter of its radii, and the additional transverse offset is at most $c'H|n-n'|$. Evaluation on this segment and the same norm calculation give the last bound in \eqref{eq:cylinder-comparison}.
\end{proof}

\subsection{The normalized limit}

For $U_i$ at $(z_i,r_i)$, choose a fitted slope $P_i$ for the good parameter $\Lambda\lambda_i$ of Lemma~\ref{lem:bad-scales}. Put $c_i=|P_i|$. By \eqref{eq:physical-slope-bounds}, after discarding finitely many indices,
\begin{equation}\label{eq:base-norm-bounds}
  \tfrac12\le c_i\le L+1.
\end{equation}
Use orthonormal coordinates with final unit vector $P_i/c_i$, write the transverse variable as $x\in\R^d$, and define $v_i$ by
\begin{equation}\label{eq:anisotropic-rescaling}
  U_i\bigl(z_i+r_i(\sqrt{\lambda_i}x,t)\bigr)
   =U_i(z_i)+r_i c_i\bigl(t+\lambda_i v_i(x,t)\bigr).
\end{equation}
These coordinates depend on $i$. The domains of $v_i$ exhaust $\R^d\times\R$, since $D(z_i)/r_i\ge1/s_i\to\infty$.

For a smooth test $\phi$ for $v_i$, the corresponding test for $U_i$ has gradient $c_i(\sqrt{\lambda_i}\nabla_x\phi,1+\lambda_i\phi_t)$ and Hessian
\[
  \frac{c_i}{r_i}
  \begin{pmatrix}
    D^2_{xx}\phi&\sqrt{\lambda_i}D_{xt}\phi\\
    \sqrt{\lambda_i}D_{tx}\phi&\lambda_i\phi_{tt}
  \end{pmatrix}.
\]
Consequently its infinity Laplacian is the positive factor $\lambda_i c_i^3/r_i$ times
\begin{equation}\label{eq:rescaled-pde}
  D^2\phi[(\nabla_x\phi,1+\lambda_i\phi_t)].
\end{equation}
Thus $v_i$ solve the equation obtained by setting \eqref{eq:rescaled-pde} equal to zero in the viscosity sense.

\begin{proposition}\label{prop:nonaffine-limit}
For the sequence supplied by Lemma~\ref{lem:bad-scales} and the normalization \eqref{eq:anisotropic-rescaling}, a subsequence of the $v_i$ converges locally uniformly to a nonaffine function $v\in C(\R^d\times\R)$ satisfying \eqref{eq:limit-pde} and \eqref{eq:limit-fit} with one finite constant $C_0$.
\end{proposition}

\begin{proof}
Set $\theta_i=\sqrt{\Lambda\lambda_i}$. A good fit of radius $Rr_i$ has error $\theta_i^2Rr_i$ and transverse radius $\theta_i Rr_i$. Its slope obeys uniform upper and lower norm bounds by \eqref{eq:physical-slope-bounds}, once $i$ is large enough.

For each fixed $R\ge2$, Lemma~\ref{lem:bad-scales}, applied with any fixed $M>R$, provides fits at $z_i$ through successive doublings from $r_i$, followed if necessary by a final step to $Rr_i$ of ratio at most two, for all sufficiently large $i$. Comparing them by Lemma~\ref{lem:cylinder-comparison} gives, for a slope $P_R$ at the latter radius and its unit vector $n_R$,
\begin{equation}\label{eq:base-cylinder-comparison}
  |n_R-P_i/c_i|\le C\theta_i(1+\log R),\qquad
  \bigl||P_R|-c_i\bigr|\le C\theta_i^2(1+\log R).
\end{equation}
Fix a box $|x|,|t|\le K$ in the coordinates of \eqref{eq:anisotropic-rescaling} and a radius factor $S>0$. Choose $R\ge2$ so large that $R>16K$ and
\[
  K\bigl(\Lambda^{-1/2}+C(1+\log R)\bigr)<R/4.
\]
Such a choice is possible because $R/\log R\to\infty$. Next fix $M>\max\{2,2K,R,S\}$. The physical centers in the box are within $2Kr_i$ of $z_i$, and all radii in the comparisons below are at most $\max\{R,S\}r_i$. Lemma~\ref{lem:bad-scales} therefore supplies these fits uniformly once $i$ is sufficiently large. The physical image of the box then lies inside the cylinder at $z_i$ for $P_R$, with axial and transverse offsets at most one quarter of the respective radii. For the axial offset use $|\sqrt{\lambda_i}x|+|t|\le2K$. For the transverse offset, \eqref{eq:base-cylinder-comparison} gives the bound
\[
  K\bigl(\sqrt{\lambda_i}+C\theta_i(1+\log R)\bigr)<\theta_i R/4
\]
in units of $r_i$.

Compare a fit centered at any point of this box, first at radius $Rr_i$ with the enclosing fit, and then through successive doublings or halvings, with a final ratio at most two, to radius $Sr_i$. The number of comparisons depends on $K,S$ but not on $i$. If $P=(P_x,P_t)$ is any such slope in the coordinates of \eqref{eq:anisotropic-rescaling}, Lemma~\ref{lem:cylinder-comparison} gives
\[
  |P/|P|-P_i/c_i|\le C_{K,S}\theta_i,
  \qquad \bigl||P|-c_i\bigr|\le C_{K,S}\theta_i^2.
\]
Since $P_i/c_i$ is the last coordinate vector,
\[
  |P_x|\le C_{K,S}\theta_i,\qquad
  |P_t-c_i|\le\bigl||P|-c_i\bigr|
      +\tfrac12|P|\,|P/|P|-P_i/c_i|^2
      \le C_{K,S}\theta_i^2.
\]
It follows that the normalized coefficients
\begin{equation}\label{eq:normalized-fit-slopes}
  p_i'=\frac{P_x}{c_i\sqrt{\lambda_i}},\qquad
  q_i'=\frac{P_t/c_i-1}{\lambda_i}
  \quad\text{satisfy}\quad |p_i'|+|q_i'|\le C_{K,S}.
\end{equation}
In particular, using the enclosing fit at $z_i$ in \eqref{eq:anisotropic-rescaling} gives on the box
\[
  |v_i(x,t)-p_{i,R}'\cdot x-q_{i,R}'t|\le\Lambda R/c_i.
\]
Thus $v_i$ are uniformly bounded on each fixed box.

We next express the domain of a fit in these coordinates. Suppress the index $i$ temporarily, write $\lambda=\lambda_i$, and set
\[
  A=1+\lambda q_i',\qquad N=(A^2+\lambda|p_i'|^2)^{1/2},
  \qquad n=\frac{(\sqrt\lambda p_i',A)}{N}.
\]
For fixed $K,S$, \eqref{eq:normalized-fit-slopes} gives $A>0$, $N/A=1+O_{K,S}(\lambda)$, and $|n_x|=O_{K,S}(\sqrt\lambda)$ for large $i$. The physical axis has ratio $\sqrt\lambda\,p_i'/A$ of transverse to axial displacement, which becomes the finite shear $p_i'/A$ after the transverse stretch. Figure~\ref{fig:anisotropic-fit} illustrates this change of geometry. The next calculation identifies a sheared box contained in the cylinder's image in normalized coordinates, with a fixed margin at the cylinder's axial ends. The exact decomposition
\begin{equation}\label{eq:cylinder-shear-decomposition}
  (\sqrt\lambda h,\ell)
   =\frac{\ell N}{A}\,n
       +\left(\sqrt\lambda\left(h-\frac{p_i'}A\ell\right),0\right)
\end{equation}
shows that the mapped increment lies in the radius-$Sr_i$ cylinder whenever
\begin{equation}\label{eq:normalized-fit-domain}
  |\ell|\le S/2,\qquad \left|h-\frac{p_i'}A\ell\right|\le S,
\end{equation}
for all large $i$. Indeed the residual in \eqref{eq:cylinder-shear-decomposition} has transverse projection at most $\sqrt\lambda S\le\theta_i S$, and its axial projection is $O_{K,S}(\lambda S)$. The first term has axial length at most $(1+O_{K,S}(\lambda))S/2$, leaving a strict axial margin. Dividing \eqref{eq:physical-fit} by $r_i c_i\lambda_i$ therefore gives
\begin{equation}\label{eq:normalized-fit}
  |v_i(x+h,t+\ell)-v_i(x,t)-p_i'\cdot h-q_i'\ell|
      \le\frac{\Lambda}{c_i}S
\end{equation}
on \eqref{eq:normalized-fit-domain}. The coefficient of $S$ is independent of $K$ and $S$.

\begin{figure}[tbp]
\centering
\begin{tikzpicture}[x=1cm,y=1cm,>=Latex,font=\small]
  \definecolor{fitblue}{RGB}{35,88,136}
  \definecolor{fitgreen}{RGB}{46,112,85}
  \node at (2.25,4.68) {physical cylinder};
  \node at (10.25,4.68) {limiting sheared box};
  \node[anchor=north,align=center] at (2.25,-0.76)
      {$Z_\lambda(z,r;P)$};
  \begin{scope}[shift={(2.25,1.85)},rotate=-16]
    \filldraw[draw=fitblue,fill=fitblue!9,line width=0.7pt]
        (-0.39,-1.50) rectangle (0.39,1.50);
    \draw[fitblue,densely dashed] (0,-1.67)--(0,0);
    \draw[fitblue,->,line width=0.85pt] (0,0)--(0,1.96)
        node[above right=1pt] {$P$};
    \fill (0,0) circle (2pt);
    \node[below left=3pt] at (0,0) {$z$};
    \draw[<->] (-0.74,0)--(-0.74,1.50)
        node[midway,left=4pt] {$r$};
    \draw[black!50] (-0.80,0)--(-0.48,0)
                        (-0.80,1.50)--(-0.48,1.50);
    \draw[<->] (0,-1.81)--(0.39,-1.81);
    \draw[black!50] (0,-1.70)--(0,-1.89)
                        (0.39,-1.60)--(0.39,-1.89);
    \node[below=4pt] at (0.28,-1.81) {$\sqrt\lambda\,r$};
  \end{scope}
  \draw[black!45,dotted] (2.25,0.10)--(2.25,4.05);
  \fill (2.25,1.85) circle (2pt);
  \node[anchor=east,font=\footnotesize,black!70] at (1.58,3.93)
      {base direction};
  \draw[black!45] (1.65,3.93)--(2.20,3.93);
  \draw[->,line width=0.8pt] (4.30,1.90)--(6.93,1.90);
  \node[align=center] at (5.61,2.58)
      {stretch transverse\\directions};
  \node[align=center,font=\footnotesize] at (5.61,1.26)
      {take fixed margins,\\then pass to the limit};
  \begin{scope}[shift={(10.25,1.85)}]
    \filldraw[draw=fitgreen,fill=fitgreen!9,line width=0.7pt]
        (-1.95,-1.30)--(0.65,-1.30)--(1.95,1.30)--(-0.65,1.30)--cycle;
    \draw[->] (-2.28,0)--(2.34,0) node[right] {$h$};
    \draw[->] (0,-1.57)--(0,1.89) node[above] {$\ell$};
    \draw[fitgreen,densely dashed,line width=0.85pt]
        (-0.65,-1.30)--(0.65,1.30);
    \node[fitgreen,anchor=west] at (0.98,0.57) {$h=p\ell$};
    \draw[fitgreen!70] (0.91,0.60)--(0.42,0.84);
    \fill (0,0) circle (2pt);
    \node[above left=3pt] at (0,0) {$0$};
    \draw[<->] (0.65,1.56)--(1.95,1.56)
        node[midway,above=2pt] {$R$};
    \draw[fitgreen!50] (0.65,1.35)--(0.65,1.65)
                          (1.95,1.35)--(1.95,1.65);
    \draw[<->] (-2.25,0)--(-2.25,1.30)
        node[midway,left=3pt] {$R$};
    \draw[fitgreen!50] (-2.31,1.30)--(-0.75,1.30);
    \node[align=center] at (0,-2.69)
        {$|\ell|\le R,\quad |h-p\ell|\le R$};
  \end{scope}
\end{tikzpicture}
\caption{From thin cylinders to sheared fitting domains. The left panel
 shows a two-dimensional section of $Z_\lambda(z,r;P)$, with axial radius
 $r$ and transverse radius $\sqrt\lambda\,r$; the dotted line indicates the
 base direction used for rescaling. Stretching transverse directions makes
 the small inclination of a nearby fitted axis visible as a finite shear.
 Fixed margins give contained sheared boxes, and passage to the limit
 yields $Q_R(z;p)$, shown on the right in increment coordinates $(h,\ell)$.
 The panels are schematic: the finite-scale cylinder image is not asserted
 to equal the limiting box.}
\label{fig:anisotropic-fit}
\end{figure}


This proves equicontinuity on compact sets. Given an accuracy $\varepsilon>0$, first fix $S>0$ so small that $2\Lambda S<\varepsilon/3$. For a fixed box, \eqref{eq:normalized-fit-slopes} now bounds its coefficients by some finite $C_{K,S}$. Choose the increments sufficiently small, depending on this bound and $S$, that \eqref{eq:normalized-fit-domain} holds and the affine increment in \eqref{eq:normalized-fit} is less than $\varepsilon/3$. For all sufficiently large $i$, \eqref{eq:base-norm-bounds} and \eqref{eq:normalized-fit} then bound the value increment by $2\varepsilon/3$. Discard an initial segment so that all domains contain the compact box. The threshold in $i$ may depend on $\varepsilon$; the finitely many remaining preceding continuous functions can be included by decreasing the increment size. The local bounds and domain exhaustion now give a locally uniformly convergent subsequence, with limit $v$ and $v(0,0)=0$.

The jet operators $X[(p_x,1+\lambda_i p_t)]$ converge locally uniformly to $X[(p_x,1)]$. Viscosity stability gives \eqref{eq:limit-pde}; equivalently, strictify a smooth upper or lower test for $v$, choose nearby touching points for $v_i$, and pass to the limit in \eqref{eq:rescaled-pde}. To pass the fits, fix a center and a target radius $R>0$, and take $S=4R$ in \eqref{eq:normalized-fit}. After a further subsequence its bounded coefficients converge to $p,q$. The shears $p_i'/(1+\lambda_i q_i')$ converge to $p$. Every increment with $|\ell|\le R$ and $|h-p\ell|\le R$ then satisfies \eqref{eq:normalized-fit-domain}, uniformly, for all sufficiently large $i$. Passing to the limit in \eqref{eq:normalized-fit} yields \eqref{eq:limit-fit} with
\[
  a=q+|p|^2/2,\qquad C_0=8\Lambda=256.
\]
The additional subsequence may depend on the chosen center and radius: the limit function $v$ stays fixed, and \eqref{eq:limit-fit} asserts existence separately for each choice.

Finally suppose that $v$ is affine. Since $v(0,0)=0$, write $v(x,t)=p\cdot x+qt$. In the physical coordinates consider
\[
  \widetilde P_i=c_i(\sqrt{\lambda_i}p,1+\lambda_i q).
\]
These slopes are nonzero for large $i$. The cylinders $Z_{\lambda_i}(z_i,r_i;\widetilde P_i)$ lie in one fixed box of the rescaled coordinates. In fact, if
\[
  N_i=\bigl(\lambda_i|p|^2+(1+\lambda_i q)^2\bigr)^{1/2}\longrightarrow1,
\]
a point in such a cylinder, in units of $r_i$, is $s(\sqrt{\lambda_i}p,1+\lambda_i q)/N_i+\sqrt{\lambda_i}\xi$ with $|s|\le1$, $|\xi|\le1$, and $\xi$ perpendicular to the axis. Its rescaled coordinates satisfy $|x|\le |p|/N_i+1$ and $|t|\le1+\sqrt{\lambda_i}$. On this fixed box the residual for the proposed fit is exactly
\[
  r_i c_i\lambda_i\bigl(v_i(x,t)-p\cdot x-qt\bigr)=o(\lambda_i r_i)
\]
uniformly, by local uniform convergence and the upper bound for $c_i$. It is eventually at most $\lambda_i r_i$ on the whole cylinder. This makes $\lambda_i$ good for $U_i$ at $(z_i,r_i)$, contrary to Lemma~\ref{lem:bad-scales}. Hence $v$ is nonaffine.
\end{proof}

\begin{proof}[Proof of Proposition~\ref{prop:power-decay}]
If no power bound held, Lemma~\ref{lem:bad-scales} and Proposition~\ref{prop:nonaffine-limit} would produce a nonaffine entire solution satisfying all hypotheses of Theorem~\ref{thm:liouville}, with $m=d$ and $C_0=256$. The theorem makes that solution affine, a contradiction. Hence a power bound holds. Reducing its exponent if necessary gives $\gamma\le1$. Since the family $\mathcal U$ and the function $b$ were fixed solely by $d$, the constants in \eqref{eq:power-decay} depend only on $d$.
\end{proof}

\subsection{From cylinder fits to a H\"older gradient}

We finish the reduction by extracting the uniform estimate from \eqref{eq:power-decay}. The exponent comes from the inscribed balls: their radii are $\sqrt\lambda r$, whereas the fitting error is $\lambda r$.

\begin{proof}[Proof of Theorem~\ref{thm:main}, assuming Theorem~\ref{thm:liouville}]
Take $B,\gamma,s_0$ from Proposition~\ref{prop:power-decay}. At every center $z\in\overline B_{1/2}^n$ one has $D(z)\ge d_0:=1/4$. Choose $A_0>Bd_0^{-\gamma}$. There is $r_0>0$, depending only on $d$, such that for all $0<r\le r_0$,
\[
  r/d_0\le\min\{s_0,1/8\},\qquad A_0r^\gamma<1.
\]
For every $U\in\mathcal U$ and each such center and radius,
\[
  \beta_U(z,r)\le b(r/D(z))
       \le B(r/d_0)^\gamma<A_0r^\gamma.
\]
Thus $\lambda=A_0r^\gamma$ is good, with the same constants for all functions and centers.

Fix $U\in\mathcal U$. Set
\[
  r_j=2^{-j}r_0,\qquad
  \rho_j=\sqrt{A_0}\,r_j^{1+\gamma/2},\qquad j\ge0,
\]
and at each center $z\in\overline B_{1/2}^n$ choose a fitted slope $P_j(z)$ for the parameter $A_0r_j^\gamma$. The fit holds on the inscribed ball $\overline B_{\rho_j}(z)$, with error at most $A_0r_j^{1+\gamma}$. Comparing consecutive fits on their common ball $\overline B_{\rho_{j+1}}(z)$ gives
\[
  |P_j(z)-P_{j+1}(z)|
   \le\frac{A_0(r_j^{1+\gamma}+r_{j+1}^{1+\gamma})}{\rho_{j+1}}
   \le C r_j^{\gamma/2}.
\]
The geometric tail converges uniformly in $U$ and $z$. Consequently there is a vector $P_\infty(z)$ with
\begin{equation}\label{eq:slope-tail}
  |P_j(z)-P_\infty(z)|\le C r_j^{\gamma/2}.
\end{equation}
If $0<|h|\le\rho_0$, choose $j$ so that $\rho_{j+1}<|h|\le\rho_j$. The fit at radius $r_j$ and \eqref{eq:slope-tail} give
\begin{align*}
  \frac{|U(z+h)-U(z)-P_\infty(z)\cdot h|}{|h|}
   &\le\frac{A_0r_j^{1+\gamma}}{\rho_{j+1}}
            +|P_j(z)-P_\infty(z)|\\
   &\le C r_j^{\gamma/2}\le C'|h|^{\alpha_d},
\end{align*}
where
\begin{equation}\label{eq:holder-exponent}
  \alpha_d=\frac{\gamma}{2+\gamma}\in(0,1/3].
\end{equation}
The final inequality follows from
$\rho_{j+1}=\sqrt{A_0}\,2^{-(1+\gamma/2)}r_j^{1+\gamma/2}$ and
$(1+\gamma/2)\alpha_d=\gamma/2$. Thus $U$ is differentiable at every such $z$, with gradient $P_\infty(z)$, and
\begin{equation}\label{eq:uniform-taylor-estimate}
  |U(z+h)-U(z)-\nabla U(z)\cdot h|
       \le C|h|^{1+\alpha_d}\qquad(|h|\le\rho_0).
\end{equation}
All increments here remain in $G$, since each lies in one of the admissible inscribed balls. The original Lipschitz bound gives $|\nabla U(z)|\le L$.

To compare gradients, take $z,z'\in\overline B_{1/2}^n$ with $\ell=|z-z'|\le\rho_0/2$ and $\ell>0$, and fix any unit vector $e\in\R^n$. The evaluation point $z'+\ell e$ need not lie in $B_{1/2}^n$; the Taylor estimates centered at $z$ and $z'$ still apply. Applying \eqref{eq:uniform-taylor-estimate} at $z'$ to the increment $\ell e$ gives
\[
  U(z'+\ell e)-U(z')
       =\ell\nabla U(z')\cdot e+O(\ell^{1+\alpha_d}).
\]
Applying the same estimate at $z$ to the increments $z'-z$ and $z'+\ell e-z$, whose lengths are at most $2\ell$, gives
\[
  U(z'+\ell e)-U(z')
       =\ell\nabla U(z)\cdot e+O(\ell^{1+\alpha_d}).
\]
The constants are independent of $e,U,z,z'$. Subtracting and taking the supremum over unit vectors proves
\[
  |\nabla U(z)-\nabla U(z')|\le C\ell^{\alpha_d}.
\]
For larger distances the same estimate follows, after increasing $C$, from $|\nabla U|\le L$. Restriction to points $(y,0)$ therefore gives the claimed uniform gradient norm and H\"older seminorm for every normalized $u$.

Finally, if $u$ is as in Theorem~\ref{thm:main} and its oscillation is $M>0$, then $(u-u(0))/M$ is a normalized solution: positive scaling and addition of constants preserve the viscosity equation by transformation of tests. Multiplying the normalized estimates by $M$ gives the theorem. If $M=0$, the function is constant and the conclusion is immediate.
\end{proof}

It remains to prove Theorem~\ref{thm:liouville}. The following sections use only its equation and fitting hypothesis; the original infinity-harmonic family will no longer enter the argument.
